% !TeX root = ../抽象代数笔记.tex
% 由主文档通过 \input 载入；请编译主文档。
\setcounter{chapter}{5} % 保留原书章号。
\chapter{群作用与计数}\label{ch:actions}

\begin{description}
  \item[学习目标] 把群元素看作集合上的变换，以轨道与稳定子组织计数和结构证明。
  \item[先修知识] 第\ref{ch:quotient-groups}章。
\end{description}

\section{群作用、轨道与稳定子}
定义群 $G$ 在集合 $X$ 上的左作用（group action），以及忠实作用、传递作用。对 $x\in X$ 定义轨道 $Gx$ 与稳定子 $G_x$，证明轨道划分 $X$，并构造轨道与左陪集集合之间的双射；当 $G$ 有限时得到轨道--稳定子公式 $|Gx|=[G:G_x]$。

\section{置换表示与凯莱定理}
说明一个群作用等价于一个从 $G$ 到 $X$ 的全体置换所成群的同态。利用群对自身的左乘作用证明凯莱定理（Cayley's theorem）：每个群都同构于某个置换群的子群。比较左乘作用、对子群陪集的作用与几何对称作用，分别计算其核。

\section{共轭作用与类方程}
定义共轭类（conjugacy class）、中心 $Z(G)$ 和元素的中心化子 $C_G(g)$。对有限群写出类方程，并由共轭作用解释共轭类大小等于 $[G:C_G(g)]$。证明非平凡有限 $p$ 群的中心非平凡，其中 $p$ 为素数、$p$ 群指阶为 $p$ 的幂的群；进一步证明 $p^2$ 阶群必为阿贝尔群。

\section{伯恩赛德引理与计数应用}
对有限群 $G$ 作用于有限集合 $X$，令 $X^g=\{x\in X:g\cdot x=x\}$，以双重计数证明伯恩赛德引理（Burnside's lemma）：轨道数为 $|G|^{-1}\sum_{g\in G}|X^g|$。用正方形顶点染色说明“按旋转视为相同”和“按全部对称视为相同”对应不同的作用群。

\section{例题、练习与自检}
计算 $S_3$ 的共轭类与类方程，求正方形顶点二色染色在不同等价约定下的数量。练习要求从文字描述中写出作用集合与作用群，判断作用是否忠实或传递，并通过稳定子计算轨道，避免直接把对象总数除以群阶。

